\documentclass[11pt,letterpaper]{article}
\usepackage[T1]{fontenc}
\usepackage[utf8]{inputenc}
\usepackage{lmodern,microtype}
\usepackage[margin=1in]{geometry}
\usepackage{amsmath,amssymb,amsthm,mathtools}
\usepackage{booktabs,longtable,array,enumitem,xcolor,listings,fancyhdr}
\usepackage{needspace}
\usepackage[breaklinks,colorlinks,linkcolor=blue!45!black,urlcolor=blue!45!black,citecolor=blue!45!black]{hyperref}
\usepackage{url}
\urlstyle{tt}
\setlength{\emergencystretch}{3em}
\setlength{\parskip}{0.35em}
\setlength{\parindent}{0pt}
\setlist{itemsep=0.2em,topsep=0.4em}
\definecolor{codebg}{RGB}{247,248,250}
\lstset{basicstyle=\ttfamily\footnotesize,backgroundcolor=\color{codebg},
  breaklines=true,columns=fullflexible,keepspaces=true,showstringspaces=false,
  frame=single,rulecolor=\color{black!12},xleftmargin=0.5em,xrightmargin=0.5em,
  aboveskip=0.6em,belowskip=0.6em}
\pagestyle{fancy}
\fancyhf{}
\lhead{\small K215 proof simplification}
\rhead{\small Checked component proposal}
\cfoot{\thepage}
\setlength{\headheight}{14pt}
\newtheorem{theorem}{Theorem}[section]
\newtheorem{lemma}[theorem]{Lemma}
\newtheorem{proposition}[theorem]{Proposition}
\newcommand{\C}{\mathcal C_{222}}
\newcommand{\Sat}{\operatorname{Sat}}
\newcommand{\enc}{\operatorname{enc}}
\newcommand{\thh}{\operatorname{th}}
\newcommand{\Free}{\operatorname{Free}}
\newcommand{\Seed}{\operatorname{Seed}}
\newcommand{\Req}{\operatorname{Req}}
\newcommand{\code}[1]{\texttt{\detokenize{#1}}}
\newcommand{\LeanJointModules}{116}
\newcommand{\LeanLegacyWarnings}{18}


\title{Simplifying the K215 Correctness Proofs\\[0.35em]
  \large A shared implication basis, direct seed semantics,\\
  and one Rocq enumerator-completeness argument}
\author{ProveIt proof-refactoring study}
\date{October 21, 2026}

\begin{document}
\maketitle
\begin{abstract}
The maintained K215 development already contains substantial proof compression.
Further simplification should therefore target repeated mathematical arguments and
unnecessary intermediate constructions. This report supplies three compatible
component proposals against the exact main-branch revision
\code{57541fe0cfe3af62b4cb58747d8d1d7c16b338d5}.
First, an existing six-goal implication derivation is integrated with the maintained
producer, replacing the separate unrestricted derivation: the two Lean proof files
shrink from 811 lines and 42,454 bytes to 401 lines and 18,243 bytes.
Second, a previously proposed direct named-formula seed construction is implemented
and checked. It removes the preliminary indiscernibility quotient from the Lean
endpoint's dependency path. Third, a conservative Rocq refactor shares candidate
progress, run-prefix transport, and boundary completeness while retaining every old
theorem interface. The archive includes patches, source snapshots, compiler receipts,
independent source reviews, and reproduction scripts. These are checked component
and generic-endpoint results; the complete machine-specific Lean and Rocq endpoints
have not been rebuilt for this packet.
\end{abstract}

\section{Recommendation and measured result}

Adopt proof sharing at the logical interface first, then shorten the lower semantic
route by constructing the full seed directly on the original model. The strongest
newly checked mathematical change is the latter: the completeness theorem already
supplies satisfaction of the original named formulas, including literal equality.
There is no need to pass through an equality-free model and an indiscernibility
quotient merely to recover those formulas' seed instances.

The shared implication idea and the direct-seed outline were already present in
the repository's research material. This packet's contribution is current-source
implementation, maintained-producer synchronization, explicit compatibility checks,
and fresh kernel checking. It does not claim a new smallest Turing machine or a new
logical axiom basis. The provenance distinctions matter because the most recent
merges deliberately preserved the checked production files while retaining later
proposals as inert inputs \cite{checkpoint,merge,seedproposal}.

\begin{table}[ht]
\centering\small
\begin{tabular}{@{}p{2.42in}rrp{1.15in}@{}}
\toprule
Scope & Before & Proposed & Meaning \\
\midrule
Two Lean implication files, lines & 811 & 401 & 410 fewer \\
Two Lean implication files, bytes & 42,454 & 18,243 & 57.0\% fewer \\
Those files plus producer, bytes & 52,651 & 28,760 & 23,891 fewer \\
Direct-seed endpoint closure, modules & 179 & 178 & One fewer net \\
Direct-seed closure, source bytes & 2,129,439 & 2,121,905 & 7,534 fewer \\
Both Lean changes, closure bytes & 2,129,439 & 2,097,694 & 31,745 fewer \\
Two Rocq enumerator files, lines & 410 & 403 & Old APIs kept \\
Two Rocq enumerator files, bytes & 18,100 & 17,833 & 267 fewer \\
\bottomrule
\end{tabular}
\caption{Different scopes are reported separately. Dependency-closure reductions
are static source measurements, not elapsed-time or proof-term-size measurements.}
\end{table}

The semantic refactor keeps the old quotient modules available to legacy clients.
Consequently, its six-file maintenance set grows by 6,795 bytes even though the K215
endpoint imports 7,534 fewer source bytes. This is a shorter endpoint argument and
dependency boundary, not a claim that every repository-wide size measure decreases.
The two implication-file saving and semantic closure saving are additive within the
selected endpoint closure. The various historical alternative implication proposals
are not additive.

\section{The exact baseline and preservation contract}

\subsection{Source files, not commit subjects, determine the baseline}

The review is pinned to the full commit hash given in the abstract. Its maintained
Lean sources are the phase-twelve checkpoint: 179 project modules, including 167
BusyBeaver modules, and 2,129,439 selected source bytes. Independently recomputing
the source bindings finds all 199 of the checkpoint's primary and compatibility
source hashes unchanged. This source census is not a new full build.

In particular, the checked \code{DominatedImplication.lean} has 287 lines and
13,723 bytes, while \code{Lukasiewicz.lean} has 524 lines and 28,731 bytes.
The 300-line core and 101-line unrestricted adapter from branch revision
\code{49e99c5b7cc6c2e164331961de28bad3904361fd} are preserved under the
phase-twelve upstream-merge inputs. They are not the production files at the pin.
The merge record itself explains that producer validation and Lean integration
remain separate acceptance steps \cite{merge}.

The supplied audit script reproduces the import closure, source hashes, and
public-declaration comparison. The proposed shared pair affects 33 modules
transitively in the original endpoint closure. The expensive site batches and
source-correspondence certificate modules are outside that affected set.

\subsection{What is preserved}

The literal 215-state transition table, blank-tape convention, source program,
enumerator transition function, numerical formula encoding, and numerical calculus
are unchanged. The existing endpoint is of the form
\[
 \operatorname{BlankHalts}(M_{215})
 \quad\Longleftrightarrow\quad
 \operatorname{BProv}(\mathrm{ZFCax}_s,[],\bot).
\]
Negating both sides yields the same nonhalting/consistency equivalence. The
intermediate equivalence between numerical derivability of \code{Code254.F} and
ZF refutability, and the existing ZF--ZFC refutation bridge, retain their types.

As the maintained checkpoint states, this is an external proof-assistant theorem.
The packet does not provide a new arithmetized derivation inside ZFC or PA. It also
does not replace the machine-specific validation evidence with a claim of a fresh
full endpoint build. The actual acceptance scope is given in Section~\ref{sec:checks}.

\section{One implication derivation for both logical layers}

\subsection{The abstract contract}

Let $A$ be an arbitrary carrier, let $x\to y$ be an arbitrary binary operation on
$A$, and let $\preceq$ be a relation satisfying only
\[
 u\preceq u,\qquad
 u\preceq x\Longrightarrow u\preceq(x\to y),\qquad
 u\preceq y\Longrightarrow u\preceq(x\to y).
\]
No injectivity, transitivity, semantics, truth element, or numerical representation
is required. Define
\[
 \rho_8(a,b,c,u)=((c\to u)\to a)\to((a\to c)\to(b\to c)).
\]
For a predicate $T:A\to\mathrm{Prop}$, the restricted interface supplies modus
ponens and $T(\rho_8(a,b,c,u))$ whenever $u\preceq a$.

The preserved six-goal derivation establishes unrestricted $\rho_8$, composition,
antecedent lifting, and the classical implicational basis:
\begin{align*}
 K(a,b)&=a\to(b\to a),\\
 P(a,b)&=((a\to b)\to a)\to a,\\
 S(a,b,c)&=(a\to(b\to c))\to((a\to b)\to(a\to c)).
\end{align*}
Each internal inference is an explicit modus-ponens application to already proved
instances. Each restricted-row guard is justified by reflexivity and insertion
into an implication operand, or by a local guard hypothesis. The proof is therefore
representation independent in exactly the sense required by the existing interface.

\begin{proposition}[Unrestricted adapter]
An unrestricted $\rho_8$ calculus with modus ponens is an instance of the restricted
interface with $u\preceq a$ defined to be $\mathrm{True}$.
\end{proposition}
\begin{proof}
Reflexivity and both insertion rules hold trivially. The restricted-row field
ignores its guard and invokes the supplied unrestricted row. Applying the generic
theorems gives $K$, Peirce, composition, and $S$ with the old unrestricted types.
\end{proof}

This small interpretation is the entire new unrestricted proof. It replaces
\code{cd_0} through \code{cd_106}; there is no second detachment development to
maintain. The numerical adapter remains byte-for-byte unchanged and continues to
instantiate the same generic interface with the repository's coding operation and
natural-number domination relation. Its legacy truth field remains available;
the generic theorem does not use it.

\subsection{Proof accounting and interface preservation}

The raw shared derivation has 58 private modus-ponens nodes. Two final aliases are
inlined into their public results, leaving 56 private nodes and three public
modus-ponens applications, including the lifting wrapper common to the comparison.
Thus both raw and rendered sources contain 59 explicit MP applications when that
wrapper is counted. The six goals do not add a second proof tree.

The public \code{Language}, \code{Restricted}, \code{rho8}, \code{rho8_all},
\code{composition}, and \code{lift_right} interfaces remain unchanged. The generic
module gains $K$, Peirce, and $S$ exports. The unrestricted \code{Implicational}
structure, its field order, and all four basis theorem types are preserved.
The final Tarski namespace section is byte-identical. Fourteen original public
declaration headers compare identically as normalized source text; separate typed
Lean clients check elaboration, not just text.

The core grows by 13 lines while the unrestricted module shrinks by 423 lines.
The relevant unit of comparison is both files together. The resulting 410-line
reduction is neither a minimality theorem for Hilbert proofs nor a bound on fully
expanded proof terms.

\subsection{Synchronizing the maintained producer}

Changing only the Lean files would leave a fragile integration: the maintained
producer would still pin and regenerate the older three-goal core. The delivered
patch therefore changes \code{generate_dominated_implication_lean.py} as well.
Its successor reads the raw six-goal source and derivation ancestry already
preserved under phase twelve. It does not execute the historical search program
or historical renderer.

The producer checks the pinned input hashes, replays MP and guard obligations in
the raw source, performs the two explicit alias substitutions, and replays the
inlined source with all six public targets. It checks 23 restricted-row uses in
each source form. The separate condensed-detachment text is pinned ancestry data;
it is not represented as having been parsed by that source reader. The new source
reader is a diagnostic and reproducibility check. Lean's kernel is the authority
for the exported theorems.

The output is exactly the preserved 14,146-byte core and 4,097-byte unrestricted
adapter, plus the unchanged 2,680-byte numerical adapter. The producer grows by
320 bytes. Counting the two changed proof files and the producer gives a net
reduction of 23,891 bytes and 401 lines.

The authoritative maintenance command after applying the patch is
\begin{lstlisting}
python3 Computability/BusyBeaver/Research/ZFC/k225-formal/\
  generate_dominated_implication_lean.py --check
\end{lstlisting}
Use \code{--out DIR} to regenerate into a separate Lean source root. The preserved
core's historical header names the research producer; the command above is the
maintained successor and intentionally uses the preserved six-goal inputs.

\section{Direct full seed semantics on the original domain}

\subsection{Why the old quotient can be bypassed}

The old lower chain first derives seed instances for equality-free formulas,
then forms an indiscernibility quotient, then transfers the seed property to
all named formulas on that quotient. That route is valid, but its starting
interface is stronger than the route uses.

Write $\thh(f)$ for the numerical calculus's equality interpretation of a named
formula, and set
\[
 \mathrm{DerivableNF}(f)\;:\!\Longleftrightarrow\;
 \C\bigl(\enc(\thh(f))\bigr).
\]
The syntactic premise of \code{KEind222.c222_complete} mentions the interpreted
code. Its semantic premise, however, supplies
\[
 H_T:\quad \forall f\in\mathrm{DerivableNF}\ \forall e,
                  \Sat(E,e,f).
\]
This is satisfaction of the \emph{original} named formula $f$, with literal
equality. It is not merely satisfaction of $\thh(f)$. The direct construction
uses that exact interface. It never assumes that satisfaction of an equality
translation can be reflected back on the original domain.

\subsection{The seed property and named bundle}

For a membership relation $E$ and a binary relation $R$, write
\[
 x\subseteq_E y\;:\!\Longleftrightarrow\;\forall w\,(E(w,x)\Rightarrow E(w,y))
\]
and define the request condition
\[
 \Req_R(x,C)\;:\!\Longleftrightarrow\;
 \forall y\,[R(x,y)\Rightarrow\exists z\,(R(x,z)\land E(z,C))].
\]
The existing Glazer seed condition can then be written compactly as
\begin{align*}
 \Seed_E(R,A)\;:\!\Longleftrightarrow\;\exists C\;[\,&\Req_R(A,C)\ \land\\
 &\forall y\,(E(y,C)\Rightarrow\exists x\,[R(x,y)\land
          \forall x'\,(x'\subseteq_E y\Rightarrow\Req_R(x',C))])\,].
\end{align*}
This is the repository's \code{SeedProp}; no new set-theoretic principle is
introduced by this notation.

The named formula \code{bundleNF a b k m} expresses that condition using input
name $a$, output name $b$, private container name $k$, and relation matrix $m$.
The key new semantic lemma is the named analogue of the existing numerical
bundle lemma.

\begin{lemma}[Named bundle semantics]
Suppose $a\ne b$, $k\ne a$, $k\ne b$, and $k$ is not free in $m$. Then
\[
 \Sat(E,e,\operatorname{bundleNF}(a,b,k,m))
 \quad\Longleftrightarrow\quad
 \Seed_E\bigl((x,y)\mapsto\Sat(E,e[a:=x][b:=y],m),e(a)\bigr).
\]
No irreflexivity hypothesis on $E$ is required for this lemma.
\end{lemma}
\begin{proof}
Expand the named request macro to the request condition and the closure macro
to the second conjunct of the seed condition. The auxiliary variable-relation
implication is true because $k$ differs from $a$ and $b$. Classical normalization
of a negated universal implication supplies a container $C$. Commute updates
at distinct names and use freshness to remove the update $e[k:=C]$ from the
matrix's satisfaction predicate. The two resulting conjuncts are the seed
conjuncts in reversed order; commutativity of conjunction finishes the proof.

Named falsity is actual falsity, which is why this lemma needs no irreflexivity
premise. The existing numerical bundle theorem still needs its original
irreflexivity hypothesis because its decoded falsity code has a different
representation. That numerical theorem is left unchanged.
\end{proof}

\subsection{A fresh private marker without renaming}

For an arbitrary named formula $\psi$, including one with equality, put
\[
 m=\neg\neg\psi,\qquad k=\enc(\thh(m)),\qquad
 B_\psi=\operatorname{bundleNF}(0,1,k,m).
\]
The following facts justify the construction.

\Needspace{7\baselineskip}
\begin{lemma}[Freshness and role bounds]
The marker satisfies $1<k$ and $\neg\Free(k,m)$.
\end{lemma}
\begin{proof}
First show by structural induction that
\[
 \Free(i,f)\Longrightarrow\Free(i,\thh(f)).
\]
The equality case uses the freshness of the auxiliary name in the membership
indiscernibility formula. Its two original free names remain free. The universal
case preserves the same binder. This is a free-support statement, not a claim
that every syntactic occurrence survives translation.

The existing code-bound theorem says that a name occurring in a formula is
strictly smaller than its encoding. If $k$ were free in $m$, free-support
preservation and that code bound would imply
$k<\enc(\thh(m))=k$, a contradiction.

For the role bound, translation commutes with negation. The outer negation in
$m$ is an implication to the falsity code, and the existing right-projection
bound for the pairing operation makes its encoding greater than one. The
delivered proof obtains this directly from \code{le_P_right}; it does not add
another positivity induction over $\psi$. Double negation simultaneously
preserves the relation defined by $\psi$.
\end{proof}

\begin{lemma}[The required bundle is derivable]
$B_\psi$ belongs to $\mathrm{DerivableNF}$.
\end{lemma}
\begin{proof}
Every atom introduced by the named bundle constructors is a membership atom.
Consequently equality translation commutes with those constructors:
\[
 \thh(\operatorname{bundleNF}(a,b,k,m))
  =\operatorname{bundleNF}(a,b,k,\thh(m)).
\]
With $k=\enc(\thh(m))$, the existing encoding identity gives
\[
 \enc(\thh(B_\psi))=\operatorname{bundle222}(0,1,k).
\]
Since $0<k$ and $1<k$, this is an enabled \code{C222.bundle} instance. By the
definition of $\mathrm{DerivableNF}$, the claim follows.
\end{proof}

\begin{theorem}[Direct full seed]
Let $E$ be irreflexive and suppose $H_T$ gives satisfaction of every formula in
$\mathrm{DerivableNF}$ under every valuation. Then $E$ has the seed property for
every named relation on its original domain, including relations defined using
literal equality.
\end{theorem}
\begin{proof}
For $\psi$ and a valuation $e$, apply $H_T$ to $B_\psi$. The preceding lemmas
allow application of named bundle semantics, producing the seed property for
the relation defined by $m=\neg\neg\psi$. Classical double-negation elimination
gives pointwise equivalence with the relation defined by $\psi$. The existing
\code{seedProp_congr} transports the seed property across that equivalence.
Combine these instances with the supplied irreflexivity hypothesis to obtain
\code{FullSeed mem}.
\end{proof}

\subsection{The actual endpoint change}

The lower endpoint now uses the original domain and its existing valuation:
\begin{lstlisting}
have hfull := fullSeed_direct_of_derivable mem hirr hT
exact zf_consistent_of_formSeed (formSeed_of_fullSeed hfull)
  hfull.irr (e 0) hf
\end{lstlisting}
The witness \code{e 0} replaces the old quotient representative
\code{KIndisc222.mk mem (e 0)}. Thus no inhabited-domain premise is dropped.
The full lower theorem's statement is unchanged.

To make the dependency reduction real, \code{KSeedInterface222} holds the named
bundle syntax, its encoding identity, \code{DerivableNF}, and \code{FullSeed},
retaining their original namespaces and types. \code{KSeedDirect222} supplies the
new proof. The old equality-free and indiscernibility modules import the shared
interface and still compile as compatibility clients.

Three modules leave the selected endpoint closure:
\code{Code254Models}, \code{KSeedModel222}, and \code{KIndisc222}. Two enter:
\code{KSeedInterface222} and \code{KSeedDirect222}. The resulting reduction is
150 selected source lines and 7,534 bytes. The first quotient is bypassed; the
later tagged-domain bisimulation quotient, well-founded part, and ZF model
construction remain unchanged. The same architectural idea already appears in
the Rocq lower proof, so no corresponding second deletion is claimed there
\cite{seedproposal,rocqlower}.

\section{One Rocq enumerator-completeness argument}

\subsection{The common command boundary}

The maintained Rocq lane separately proves that enumeration reaches later
candidates from its zero start and from its bootstrap seed. It also repeats the
run-prefix transport and final completeness argument. The refactor makes the
shared hypothesis explicit: a state is at a command boundary when its pending
command field \code{pl} is zero \cite{enumerator}.

For an arbitrary boundary state $s$ and $j\in\mathbb N$, the shared reachability
theorem gives
\[
 \operatorname{halts\_from}(s)\quad\lor\quad
 \exists n,t\;[\operatorname{run}(n,s)=\operatorname{Some}(t)
       \land\mathrm{pl}(t)=0\land\mathrm{np}(t)=\mathrm{np}(s)+j].
\]
The proof is the existing generalized seed induction moved to the common compiler
module. It is complemented by one shared lemma transporting halting backward
along an already proved finite run.

\begin{theorem}[Boundary completeness]
If $\mathrm{pl}(s)=0$ and \code{thm222 1}, then
\code{halts_from s}.
\end{theorem}
\begin{proof}
Let $Y=\mathrm{np}(s)$ and compile a refutation into a command block $L$.
Choose the reset-prefixed candidate code
\[
 N=P\bigl(3,P(Y,P(1,\operatorname{enc\_tail}(L\mathbin{+\!+}\mathrm{SUFFIX})))\bigr).
\]
The inherited pairing inequalities give $Y\le N$. Apply boundary reachability
with $j=N-Y$. If the search has already halted, the conclusion follows.
Otherwise it reaches a boundary state with candidate $N$. The existing
\code{program_halts} contract runs the reset-prefixed refutation from every
remaining working-register tuple. Compose that halting run with the prefix.
\end{proof}

This is a completeness theorem. It does not claim that arbitrary initial working
registers satisfy the invariant needed for soundness. That separation is why the
generalization can have the small hypothesis \code{pl seed = 0} without weakening
the established consistency endpoint.

The old zero-start and bootstrap completeness theorems become direct
specializations. All 46 original theorem headers are preserved, including the
old drain, progress, and prefix helpers. The maintained templates, generated
namespace variants, and corresponding manifest hashes are synchronized.

\subsection{Scope of the improvement}

The two touched proof files shrink by only seven lines and 267 bytes. Retaining
every old helper interface costs compatibility wrappers, so this proposal is
primarily a reduction in repeated arguments. The repository already contains
a more invasive arbitrary-state rank proposal; that proposal removes some of
those interfaces. The present component is a conservative alternative, not a
new claim of scheduler-theorem novelty \cite{rank}.

\section{Verification and the remaining acceptance boundary}\label{sec:checks}

\subsection{Completed checks}

The final joint Lean validation freshly compiles \LeanJointModules{} project
modules under Lean 4.32.0 in one clean replay into a new output directory, using
no precompiled project objects. It checks the shared implication pair, both new
seed modules, the modified generic \code{KEndpoint}, retained legacy clients,
and typed compatibility audits. The receipt records 26 selected named assumption
reports. The distributed standalone reproducer was used for this complete replay.

The new generic implication results report no axioms. The named semantic and
generic endpoint reports use subsets of the existing standard Lean axioms
\code{propext}, \code{Classical.choice}, and \code{Quot.sound}. No new axiom,
\code{sorry}, or unchecked native decision is introduced. The fresh direct-proof
and audit modules have zero warnings; \LeanLegacyWarnings{} preexisting linter
warnings occur in unchanged upstream sources. Exact source, object, and log
hashes are recorded with the named reports.

For Rocq, the baseline and candidate each freshly compile an eleven-module
component closure with Coq 8.18.0 and OCaml 4.14.1. The baseline's nine typed
interface reports and the candidate's twelve reports are closed under the
global context. Both recursive \code{coqchk} runs succeed. The build copies
change only the \code{From Stdlib Require} spelling to \code{From Coq Require}
for this toolchain; distributed production files retain their original spelling.

Source checks additionally confirm that the patches apply at the pinned revision,
the producer reproduces the intended Lean sources exactly, the numerical adapter
is unchanged, all Rocq original theorem headers remain identical, and the generated
Rocq files match the maintained templates and provenance headers.

\subsection{What these checks do not establish}

The joint Lean replay includes the generic machine-to-ZFC endpoint, not a fresh
build of \code{KEndpoint215} and its complete machine-specific closure. The Rocq
checks cover the compiler/seed component, not the complete production Rocq 9.1
machine-to-ZFC build. Existing machine and decoder receipts are not relabeled as
new checks. Import-count and source-size measurements are not runtime benchmarks.

The final adoption gate is to rebuild the affected descendants and the actual
machine-specific endpoints with the production toolchains, then compare their
exact theorem types and assumption profiles. The unchanged site and source
certificates may be reused only through the repository's established exact
source/dependency/object binding procedure. This packet leaves the older
accepted receipts intact.

\section{Integration plan}

Use an isolated branch based on the pinned revision and apply the three patches
in the supplied order. The combined patch contains the same changes and may be
used instead of the three individual patches.

\begin{enumerate}
\item Apply the shared-basis patch. Run the maintained producer's \code{--check}
mode and the focused typed API check. This changes two Lean proof files and
their maintained producer; it does not change the numerical adapter.
\item Apply the direct-seed patch. Rebuild the new interface and direct proof,
the lower-seed client, the generic endpoint, and both legacy modules. The joint
standalone replay in this archive already verifies this combination with the
shared implication pair.
\item Apply the Rocq enumerator patch. Regenerate the two namespace variants
through the normal production pipeline and check the manifest bindings. The
archive's component replay checks both baseline and proposed interfaces.
\item Run the production machine-specific endpoint checks and preserve their
new source/dependency/object receipts. Only after this step should the changes
be labeled adopted at the machine endpoint.
\end{enumerate}

The archive includes source snapshots for offline component reproduction. It
does not contain compiled Lean or Rocq objects. The preserved phase-twelve inputs
used by the producer remain at their established repository paths; no historical
derivation is overwritten. Read the top-level \code{README.md} for exact commands
and \code{MANIFEST.sha256} for delivered file hashes.

\section{Further simplification questions}

\subsection{Share named and numerical bundle normalization}

The direct proof intentionally follows the already understood named-bundle
normalization before attempting a larger abstraction. The named and decoded
versions still share update-commutation and request/closure reasoning. Can one
small semantic interface capture both while keeping explicit the difference
between primitive named falsity and decoded numerical falsity? A successful
abstraction must preserve the numerical irreflexivity premise and the coincident
role cases needed by the upper chain. Merely deleting those premises would be
incorrect.

A second possible contraction is to construct \code{FormSeed} directly at named
roles 3 for input and 1 for output, matching \code{Code254GIC.names}. This could
remove the extra input-name wrapper used by \code{formSeed_of_fullSeed}. It needs
a separate compatibility split and a measured net-source comparison; it is not
implemented in this packet.

\subsection{Replace whole-site execution by reusable primitive suffixes}

The repository already has decoder-prefix and semantic-coverage proposals.
The remaining substantial task is to prove and instantiate the primitive and
return suffix contracts for every guarded outcome. The maintained site proof
covers 4,751 leaves and 6,176 outcomes. Factoring its repeated execution should
retain both decrement polarities, exact consumed paths including zero padding,
overflow behavior, the gap and spare-register conditions, and positive physical
progress for backward halting transfer. A compiled single-leaf example is not
complete coverage. This packet makes no claim to have discharged that global
suffix premise.

\subsection{A single abstract scheduler interface across proof assistants}

The new Rocq boundary theorem and the existing Lean generalized enumerator
argument suggest a common interface: a boundary predicate, candidate-rank
progress, an early-halting alternative, and an executable refutation for a
candidate above the current rank. A further refactor could derive both zero-start
and bootstrap results from this interface in both assistants. Measure source
growth from the abstraction against eliminated duplicate proofs; the present
Rocq result shows why preserving all legacy wrappers can make raw line savings
small even when the mathematical structure improves.

\subsection{Proof-DAG optimization with several exports}

The shared six-goal derivation shows that optimizing one export in isolation
misses useful sharing. Further search should minimize the union of dependencies
of all required exports, and should count generated proof files, adapters, and
the maintained producer together. Proof-search output must continue to be
replayed by the assistant kernels. No optimality conclusion follows from the
current 58-node derivation.

\subsection{A separate internal arithmetic endpoint}

If the long-term target requires an explicit internal formal derivation of the
machine/consistency equivalence in a weak arithmetic theory or in ZFC, that is
a separate development. The present refactors improve the external verified
argument and make its interfaces clearer. They do not close the internalization
obligation by themselves.

\section*{Attribution and reproducibility notes}

The machine, numerical calculus, semantic model construction, condensed-detachment
derivations, and original compiler proofs are inherited from ProveIt and the
sources credited there. In particular, the shared implication basis is preserved
branch work; the direct-seed construction implements the prior
\code{sharing-interfaces} outline and follows an architecture already used in
Rocq. The packet retains source attribution headers and the repository's MIT
No Attribution license. New compilation receipts establish the stated local
checks; historical receipts are identified as historical evidence rather than
silently promoted.

\clearpage
\begin{thebibliography}{9}
\small\raggedright
\bibitem{repo}
V. Reshetnikov and ProveIt contributors.
\emph{ProveIt}, pinned main revision \code{57541fe0cfe3}.
\url{https://github.com/VladimirReshetnikov/ProveIt/commit/57541fe0cfe3af62b4cb58747d8d1d7c16b338d5}.

\bibitem{checkpoint}
ProveIt contributors. \emph{Twelfth K215 proof-simplification checkpoint}.
Pinned repository path
\path{Computability/BusyBeaver/Research/ZFC/k225-formal/verification/k215/simplification/phase12/README.md}.
Source hashes independently checked by the supplied \code{audit_current.py}.

\bibitem{merge}
ProveIt contributors. \emph{Shared implication core and alternate semantic cut}.
Pinned repository path
\path{Computability/BusyBeaver/Research/ZFC/k225-formal/verification/k215/simplification/phase12/upstream-merge/README.md},
with \code{merge.json} and preserved inputs. Proposed proof lineage:
\url{https://github.com/VladimirReshetnikov/ProveIt/commit/49e99c5b7cc6c2e164331961de28bad3904361fd}.

\bibitem{seedproposal}
ProveIt research packet. \emph{Direct full seed: remove the first quotient, not the second}.
Pinned repository path
\path{Computability/BusyBeaver/Research/ZFC/k225-formal/verification/k215/simplification/proposals/sharing-interfaces/calculus/DIRECT_FULL_SEED.md}.
Previously an audited, uncompiled Lean proposal; implemented and checked here.

\bibitem{rocqlower}
ProveIt contributors. \emph{Complete222}, Rocq direct lower construction.
Pinned repository path
\path{Computability/BusyBeaver/Research/ZFC/k225-formal/rocq224/src/Complete222.v}.
See \code{tau_fr}, \code{th_fr}, \code{seed_code_fresh}, and \code{Q_seed}.

\bibitem{enumerator}
ProveIt contributors. \emph{L1\_CompPerm} and \emph{SeedPerm}.
Maintained templates under
\path{Computability/BusyBeaver/Research/ZFC/k225-formal/rocqperm/};
generated K215 sources under
\path{Computability/BusyBeaver/Research/ZFC/k225-formal/rocq/rocq-extension/coq/215/}.

\bibitem{rank}
ProveIt research packet. \emph{Checked transport and seed rank}.
Pinned proposal directory
\path{Computability/BusyBeaver/Research/ZFC/k225-formal/verification/k215/simplification/proposals/checked-transport/rocq_seed_rank/}.
Alternative arbitrary-state completeness refactor, credited separately.
\end{thebibliography}
\end{document}
